% Source PDF pages 100--102; manual pages 2-71--2-73.
\subsection{Radar Observations}\label{sec:radar-observations}

TAOS simulates radar observations by viewing a vehicle trajectory from a fixed
position representing a radar antenna. Any number of radar locations can be used.
For each station, a local geodetic horizon coordinate system is defined with its
origin at the antenna, as shown in \cref{fig:radar-station-coordinates}. The
station position relative to the earth center is $\vec r_r$, and its basis vectors
$\hat x_r$, $\hat y_r$, and $\hat z_r$ are computed from the station longitude and
latitude using \cref{eq:local-geodetic-x-unit-vector,eq:local-geodetic-y-unit-vector,eq:local-geodetic-z-unit-vector}.

\begin{figure}[htbp]
  \centering
  \resizebox{0.72\linewidth}{!}{\begin{tikzpicture}[>=Latex,line cap=round,line join=round,scale=0.88]
  \coordinate (c) at (-3.5,-1.9);
  \coordinate (r) at (1.0,0.55);
  \draw[very thick] (c) arc[start angle=180,end angle=0,x radius=3.5,y radius=2.25];
  \draw[->,thick] (c)--(3.8,-1.9) node[right] {equatorial plane};
  \draw[->,thick] (c)--(-3.5,2.2) node[above] {$\hat z_{\oplus}$};
  \draw[<->] (-3.8,-1.9)--(-3.8,0.45) node[midway,left] {$R_p$};
  \draw[<->] (-3.5,-2.25)--(3.1,-2.25) node[midway,below] {$R_{\oplus}$};
  \draw[->,very thick] (c)--(r) node[midway,above left] {$\vec r_r$};
  \fill (r) circle (2.3pt) node[above right] {radar station};
  \draw (0.25,0.95)--(1.65,-0.25)--(2.45,0.65)--(1.05,1.85)--cycle;
  \draw[->,thick] (r)--(0.15,1.45) node[above left] {$\hat x_r$};
  \draw[->,thick] (r)--(1.85,1.28) node[above right] {$\hat y_r$};
  \draw[->,thick] (r)--(0.55,-0.35) node[below] {$\hat z_r$};
  \node[align=left] at (2.9,0.2) {geodetic local\\horizon plane};
\end{tikzpicture}
}
  \caption{Radar station coordinates.}
  \label{fig:radar-station-coordinates}
\end{figure}

The vehicle position relative to the radar station is
\begin{equation}
  \Delta\vec r_r=\vec r-\vec r_r,
  \label{eq:radar-relative-position}
\end{equation}
and radar range is $\lVert\Delta\vec r_r\rVert$. The corresponding line-of-sight
unit vector is
\begin{equation}
  \hat u_r=\frac{\Delta\vec r_r}{\lVert\Delta\vec r_r\rVert}.
  \label{eq:radar-line-of-sight-unit-vector}
\end{equation}
The radar azimuth $\alpha_r$ and elevation $\epsilon_r$, shown in
\cref{fig:radar-observations}, are
\begin{equation}
  \tan\alpha_r
  =\frac{\Delta\vec r_r\cdot\hat y_r}
         {\Delta\vec r_r\cdot\hat x_r},
  \label{eq:radar-azimuth}
\end{equation}
\begin{equation}
  \sin\epsilon_r=-\hat u_r\cdot\hat z_r.
  \label{eq:radar-elevation}
\end{equation}
Azimuth is measured in the local geodetic horizon plane from $\hat x_r$, so
$0^\circ$ is north and $90^\circ$ is east. Elevation is positive above the local
horizon, away from the earth.

\begin{figure}[htbp]
  \centering
  \resizebox{0.70\linewidth}{!}{\begin{tikzpicture}[>=Latex,line cap=round,line join=round,scale=0.92]
  \coordinate (r) at (-3.2,0);
  \coordinate (v) at (2.9,1.25);
  \node[left,align=center] at (r) {radar\\station};
  \fill (v) circle (2.2pt) node[right] {vehicle};
  \draw[->,thick] (r)--(-2.1,1.15) node[above] {$\hat x_r$};
  \draw[->,thick] (r)--(2.7,-0.95) node[right] {$\hat y_r$};
  \draw[->,thick] (r)--(-3.2,-1.45) node[below] {$\hat z_r$};
  \draw[->,very thick] (r)--(v) node[pos=0.72,above] {$\Delta\vec r_r$};
  \draw[thin] (r)--(2.9,0.35)--(v);
  \draw[->] (-2.25,0.25) arc[start angle=18,end angle=42,radius=0.9]
    node[midway,above] {$\alpha_r$};
  \draw[->] (2.15,0.55) arc[start angle=0,end angle=48,radius=0.55]
    node[midway,right] {$\epsilon_r$};
  \draw (2.58,0.35)--(2.58,0.65)--(2.88,0.65);
\end{tikzpicture}
}
  \caption{Radar observations.}
  \label{fig:radar-observations}
\end{figure}

Differentiating the geometry gives range rate and range acceleration along
$\hat u_r$:
\begin{equation}
  \lVert\Delta\dot{\vec r}_r\rVert
  =\vec V_{\oplus}\cdot\hat u_r,
  \label{eq:radar-range-rate}
\end{equation}
\begin{equation}
  \lVert\Delta\ddot{\vec r}_r\rVert
  =\vec a_{\oplus}\cdot\hat u_r
   +\vec V_{\oplus}\cdot\dot{\hat u}_r.
  \label{eq:radar-range-acceleration}
\end{equation}
The first two derivatives of the line-of-sight unit vector are
\begin{equation}
  \dot{\hat u}_r
  =\frac{\vec V_{\oplus}
         -(\vec V_{\oplus}\cdot\hat u_r)\hat u_r}
        {\lVert\Delta\vec r_r\rVert},
  \label{eq:radar-unit-vector-rate}
\end{equation}
\begin{equation}
  \ddot{\hat u}_r
  =\frac{\vec a_{\oplus}
         -\lVert\Delta\ddot{\vec r}_r\rVert\hat u_r
         -2(\vec V_{\oplus}\cdot\hat u_r)\dot{\hat u}_r}
        {\lVert\Delta\vec r_r\rVert}.
  \label{eq:radar-unit-vector-acceleration}
\end{equation}
The azimuth derivatives are
\begin{equation}
  \dot\alpha_r
  =\frac{
   (\vec V_{\oplus}\cdot\hat y_r)(\Delta\vec r_r\cdot\hat x_r)
  -(\vec V_{\oplus}\cdot\hat x_r)(\Delta\vec r_r\cdot\hat y_r)}
  {(\Delta\vec r_r\cdot\hat x_r)^2
  +(\Delta\vec r_r\cdot\hat y_r)^2},
  \label{eq:radar-azimuth-rate}
\end{equation}
\begin{equation}
\begin{aligned}
  \ddot\alpha_r={}&
  \frac{
   (\vec a_{\oplus}\cdot\hat y_r)(\Delta\vec r_r\cdot\hat x_r)
  -(\vec a_{\oplus}\cdot\hat x_r)(\Delta\vec r_r\cdot\hat y_r)}
  {(\Delta\vec r_r\cdot\hat x_r)^2
  +(\Delta\vec r_r\cdot\hat y_r)^2}\\
  &-\frac{2\dot\alpha_r\left[
    (\vec V_{\oplus}\cdot\hat x_r)(\Delta\vec r_r\cdot\hat x_r)
   +(\vec V_{\oplus}\cdot\hat y_r)(\Delta\vec r_r\cdot\hat y_r)
   \right]}
  {(\Delta\vec r_r\cdot\hat x_r)^2
  +(\Delta\vec r_r\cdot\hat y_r)^2}.
\end{aligned}
\label{eq:radar-azimuth-acceleration}
\end{equation}
The elevation derivatives are
\begin{equation}
  \dot\epsilon_r
  =-\dot{\hat u}_r\cdot\hat z_r\sec\epsilon_r,
  \label{eq:radar-elevation-rate}
\end{equation}
\begin{equation}
  \ddot\epsilon_r
  =\left[\dot\epsilon_r^2\sin\epsilon_r
         -\ddot{\hat u}_r\cdot\hat z_r\right]\sec\epsilon_r.
  \label{eq:radar-elevation-acceleration}
\end{equation}

The radar aspect angle $\eta_r$ is the angle between the vehicle longitudinal axis
$\hat x_b$ and the line of sight from the vehicle back to the radar, $-\hat u_r$.
The meridional angle $\phi_r$ is formed by projecting that line of sight into the
vehicle $\hat y_b$--$\hat z_b$ plane, as shown in
\cref{fig:radar-aspect-meridional}.

\begin{figure}[htbp]
  \centering
  \resizebox{0.84\linewidth}{!}{\begin{tikzpicture}[>=Latex,line cap=round,line join=round,scale=0.85]
  \coordinate (a) at (-3.7,0);
  \draw (-5.2,1.0)--(a)--(-5.2,-1.0)--cycle;
  \node[left] at (a) {cg};
  \draw[->,thick] (a)--(-0.4,0) node[right] {$\hat x_b$};
  \draw[->,thick] (a)--(-3.7,-1.45) node[below] {$\hat z_b$};
  \draw[->,thick] (a)--(-1.5,-1.05) node[right] {$\hat u_r$};
  \draw[->] (-2.1,0) arc[start angle=0,end angle=-25,radius=1.55]
    node[midway,right] {$\eta_r$};

  \coordinate (b) at (2.1,0);
  \draw (b) circle (1.15);
  \node[above left] at (b) {cg};
  \draw[->,thick] (b)--(4.2,0) node[right] {$\hat y_b$};
  \draw[->,thick] (b)--(2.1,-1.45) node[below] {$\hat z_b$};
  \draw[->,thick] (b)--(3.55,-1.0) node[below right,align=left]
    {$\hat u_r$ projected\\onto $\hat y_b$--$\hat z_b$ plane};
  \draw[->] (3.05,0) arc[start angle=0,end angle=-34,radius=0.95]
    node[midway,right] {$\phi_r$};
\end{tikzpicture}
}
  \caption{Radar aspect and meridional angles.}
  \label{fig:radar-aspect-meridional}
\end{figure}

The aspect angle satisfies
\begin{equation}
  \sin\eta_r
  =\sqrt{(\hat u_r\cdot\hat y_b)^2+(\hat u_r\cdot\hat z_b)^2},
  \label{eq:radar-aspect-sine}
\end{equation}
\begin{equation}
  \cos\eta_r=-\hat u_r\cdot\hat x_b,
  \label{eq:radar-aspect-cosine}
\end{equation}
and the meridional angle is defined by
\begin{equation}
  \sin\phi_r
  =\frac{-\hat u_r\cdot\hat z_b}
         {\sqrt{(\hat u_r\cdot\hat y_b)^2+(\hat u_r\cdot\hat z_b)^2}},
  \label{eq:radar-meridional-sine}
\end{equation}
\begin{equation}
  \cos\phi_r
  =\frac{-\hat u_r\cdot\hat y_b}
         {\sqrt{(\hat u_r\cdot\hat y_b)^2+(\hat u_r\cdot\hat z_b)^2}}.
  \label{eq:radar-meridional-cosine}
\end{equation}
These calculations are implemented in \texttt{radar\_out} in the
\texttt{derivs} module.
